\begin{helper}
Let \(\eta,\varepsilon_1\in\mathbb R\) and \(n\in\mathbb N\).  Put
\[
L=\log n,\quad
k_R=(1+\eta)\sqrt{nL},\quad
K=\noderef{TwoBiteNaturalIndependenceScale}(1+\eta,n),\quad k=(K:\mathbb R),
\]
\[
t_1=\frac{\sqrt{nL}}{\log L}.
\]
Assume \(0<\eta\), \(0<\varepsilon_1\), \(0<n\), \(0<L\), \(0<\log L\),
\(0\le k\), \(k_R\le k\), \(0<t_1\),
\[
\frac{2k}{t_1}+1\le5(1+\eta)\log L,
\]
and
\[
\frac{2000(5(1+\eta))^2(\log L)^2L^3}
{\varepsilon_1(1+\eta)\sqrt{nL}}\le1.
\]
Then
\[
\noderef{RealChooseTwo}\left(\frac{2k}{t_1}+1\right)200L^3
\le\frac{\varepsilon_1}{10}k.
\]
\end{helper}
\begin{proof}
Let \(x=2k/t_1+1\) and \(A=5(1+\eta)\log L\).  Since \(0\le k\) and
\(0<t_1\), we have \(x\ge0\).  The upper bound in
\(\noderef{RealChooseTwoQuadraticBounds}\) gives
\[
\noderef{RealChooseTwo}(x)\le x^2.
\]
The displayed bound on \(x\) then gives \(x^2\le A^2\), so
\[
\noderef{RealChooseTwo}(x)200L^3\le A^2\,200L^3.
\]
The threshold assumption is equivalent to
\[
2000(5(1+\eta))^2(\log L)^2L^3\le\varepsilon_1 k_R.
\]
Since \(A^2=(5(1+\eta))^2(\log L)^2\), this implies
\[
A^2\,200L^3\le\frac{\varepsilon_1}{10}k_R.
\]
Finally \(k_R\le k\) and \(\varepsilon_1>0\), so
\[
\frac{\varepsilon_1}{10}k_R\le\frac{\varepsilon_1}{10}k.
\]
Combining the inequalities proves the claim.
\end{proof}
